\documentclass[10pt]{article}

\usepackage[utf8]{inputenc}

\usepackage[T1]{fontenc}

\usepackage{amsmath}

\usepackage{amsfonts}

\usepackage{amssymb}

\usepackage[version=4]{mhchem}

\usepackage{stmaryrd}



\begin{document}

Энергетич. уровни $\varepsilon_{i}(s)$, рассматриваемые как функции времени, называются а д и батическими термами. В условиях применимости А. т. в. $(\xi \ll 1)$ вероятности переходов между ними можно найти с помощью последовательных итераций системы уравнений для $C_{j}^{(i)}$. В частности, в первом порядке А.т. в. амплитуда перехода с начального терма $i$ в конечное состояние $j$ за все время взаимодействия определяется формулой



\[

\begin{gathered}

A_{i j}=C_{j}^{(i)}(1)=\int_{0}^{1} w_{j i}(s) \exp \left[\frac{i T}{\hbar} \int_{0}^{s} \Delta \varepsilon_{i j}(u) d u\right] d s, \\

\Delta \varepsilon_{i j}(u)=\varepsilon_{j}(u)-\varepsilon_{i}(u) .

\end{gathered}

\]



При $\xi \ll 1$ интеграл содержит быстроосциллирующую функцию и его можно вычислить с рассматриваемой точностью методом стационарной фазы, что приводит к следующему выражению для вероятности перехода:



\[

P_{i j}=\left|A_{i j}\right|^{2} \simeq \frac{2 \pi \hbar}{T}\left|w_{i j}\left(s_{0}\right)\right|^{2}\left|\frac{d}{d s} \Delta \varepsilon_{i j}(s)\right|_{s_{0}}^{-1} .

\]



Здесь $s_{0}-$ точка стационарной фазы, к-рая определяется из условия «квазипересечения" термов $\varepsilon_{i}\left(s_{0}\right)=\varepsilon_{j}\left(s_{0}\right)$. В общем случае решение этого уравнения сушествует при комплексном значении $s_{0}=s_{0}^{\prime}+i s_{0}^{\prime \prime}$, и расчет вероятности перехода требует специального рассмотрения эволюции системы вблизи точки $s_{0}$, поскольку приведенная выше функция $w_{i j}(s)$, вообще говоря, имеет особенность в этой точке. Вблизи $s_{0}$ А. т. в. необходимо модифицировать по аналогии с обычной квантовомеханич. теорией возмущений при вырождении уровней невозмушенной задачи, используя «правильную» линейную комбинацию адиабатич. функций (cM. [7]):



\[

\Phi(s) \simeq a_{1} \varphi_{1}(s)+a_{2} \varphi_{2}(s) .

\]



В этом выражении для удобства введены такие комбинации невозмущенных функций, соответствующих пересекаюшимся термам, к-рые удовлетворяют следующим условиям



\[

\left\langle\varphi_{1}^{2}\right\rangle=\left\langle\varphi_{2}^{2}\right\rangle=0 ; \quad\left\langle\varphi_{1} \varphi_{2}\right\rangle=1

\]



что всегда можно сделать при определенном выборе коэффициентов в выражениях



\[

\varphi_{1,2}(s)=C_{1,2}^{(i)}(s) \Psi_{i}(s)+C_{1,2}^{(j)}(s) \Psi_{j}(s) .

\]



Тогда уравнение Шредингера для определения термов вблизи точки $s_{0}$ сводится к системе уравнений для коэффициентов $a_{1}, a_{2}$ :



\[

\begin{gathered}

H_{11}(s) a_{1}+H_{12}(s) a_{2}=\varepsilon a_{2}, \\

H_{12}(s) a_{1}+H_{22}(s) a_{2}=\varepsilon a_{1},

\end{gathered}

\]



Tо есть



\[

\varepsilon_{1,2}(s)=H_{12} \pm \sqrt{H_{11} H_{22}} ; \quad a_{2} / a_{1}= \pm \sqrt{H_{11} / H_{12}},

\]



где



\[

H_{\alpha \beta}(s)=\left\langle\varphi_{\alpha} \hat{H} \varphi_{\beta}\right\rangle, \alpha, \beta=1,2 .

\]



Так как, по предположению, $\varepsilon_{1}\left(s_{0}\right)=\varepsilon_{2}\left(s_{0}\right)$, то один из матричных элементов (для определенности $H_{11}$ ) обращается в нуль в этой точке, причем в силу регулярности гамильтониана как функции времени при $s \rightarrow s_{0}$



\[

H_{11}(s) \simeq \operatorname{const}\left(s-s_{0}\right) \text {. }

\]



Следовательно, пересечению термов соответствует точка ветвления в комплексной плоскости переменной $s$



\[

\varepsilon_{1,2} \simeq \varepsilon\left(s_{0}\right) \pm \mathrm{const} \sqrt{s-s_{0}},

\]



причем $\varepsilon_{1}(s)$ является аналитич. продолжением функции $\varepsilon_{2}(s)$ при обходе вокруг точки $s_{0}$.



Задача вычисления вероятности перехода между термами при «квазипересечении» эквивалентна задаче о расчете надбарьерного отражения в квазиклассич. приближении (см. [7]). В частности, вероятность того, что система, находившаяся на терме $\varepsilon_{i}$ при $s=0$, перейдет на терм $\varepsilon_{j}$ при $s=1$, определяется формулой



\[

P_{i j}=\exp \left\{-\frac{2 T}{\hbar}\left|\operatorname{Im} \int_{\operatorname{Re} s_{0}}^{s_{0}}\left[\varepsilon_{i}(u)-\varepsilon_{j}(u)\right] d u\right|\right\} \equiv \exp \left\{-\frac{2 T}{\hbar} \Delta\right\},

\]



причем параметр Месси $\Delta$ вычисляется по контуру, проходяшему вдоль линии разреза, проведенного через точку $s_{0}$.



Более сложный случай $N$-уровнего пересечения термов также можно исследовать в рамках А. т. в. (см. [4]). Специального рассмотрения требует применение А. т. в. в случае, когда адиабатич. термы в нек-рой области параметров гамильтониана переходят в непрерывный спектр, как, напр., при описании процесса ионизации в медленных столкновениях атомов (см. [8]). В этом случае удобно (см. [4]) осуществить замену переменных $s$ на $\varepsilon$ с помощью функщий $s_{i}(\varepsilon)$, обратных по отношению к адиабатич. термам $\varepsilon_{i}(s)$. Тогда с помощью аналитич. продолжения адиабатич. функщий $\Psi_{i}\left[s_{i}(\varepsilon)\right]$ в комплексную плоскость переменной $\varepsilon$ удается в рамках А. т. в. вычислить вероятности переходов системы в состояния непрерывного спектра в виде разложения по параметру



\[

\xi=\frac{1}{T}\left|\frac{\partial s}{\partial \varepsilon}\right| \text {. }

\]



Лит.: [1] Эре н фест П., Относительность. Кванты. Статистика, M., 1974; [2] Born M., Fock V., "Z. Phys.», 1928, Bd 51, S. 165; [3] Born M., Oppenheimer R., "Ann. Phys.», 1927, Bd 84, S. 457; [4] Солов ве в Е. А., «Успехи физических наук», 1989, т. 157, вып. 3, с. 437; [5] Боголюбов Н.Н., Избранные труды, т. 2, Киев, $1970 ;$ [6] Л и фш и ц Е. М., П и т а е в ск и й Л. П., Статистическая физика, ч. 2, М., 1978; [7] Л а н д а у Л. Д., Л ифши ц Е. М., Квантовая механика, 3 изд., М., 1974; [8] Никити н Е. Е., Уман ск и й С.Я., Неадиабатические переходы при медленных атомных столкновениях, М., $1979 . \quad$ И. Д. Феранчук.



АДИАБАТИЧЕСКИЙ ИНВАРИАНТ - физическая величина, остающаяся практически неизменной при медленном (адиабатическом), но не обязательно малом, изменении внешних условий, в которых находится система, либо самих характеристик системы (внутреннее состояние, масса, электрический заряд и пр.). Отмеченное изменение должно происходить за времена $(\tau)$, значительно превышаюшие характерные периоды движения системы $(T)$.



В классич. механике А. и. являются переменные действия



\[

I_{k}=\oint p_{k} d q_{k}

\]



где $p_{k}-$ обобщенный импульс, $q_{k}$ - обобщенная координата, интегрирование производится по периоду (или квазипериоду).



Для гармонич. осциллятора А. и. является отношение его энергии к частоте. Характерно, что при адиабатич. изменении условий становятся связанными между собой физич. величины, к-рые вообще независимы, напр. амплитуда колебаний маятника и его длина.



Физически важным примером А. и. служит магнитный момент, создаваемый током заряженной частицы при ее движении в медленно меняющемся (в пространстве или во времени) магнитном поле: $p_{\perp}^{2} / H=\mathrm{const}$, где $p_{\perp}-$ проекция импульса заряженной частицы на плоскость, перпендикулярную направлению магнитного поля $(\boldsymbol{H})$ в данной точке пространства.



На сохранении А. и. основано так наз. дрейфовое приближение, широко используемое в физике плазмы, а также действие «магнитных пробок» и основанных на них адиабатич. ловушек - пробкотронов, применяемых в исследованиях по удержанию горячей плазмы для целей управляемого термоядерного синтеза и осуществляющихся, напр., в магнитном поле Земли.



Количество А. и. не превышает числа степеней свободы, по к-рым движение системы финитно (ограничено в пространстве). Так, в магнитных ловушках, кроме магнитного момента, может сохраняться продольный А. и., соответствующий движению вдоль магнитных силовых линий:



\[

\int_{a}^{b} p_{\|} d l

\]



где $p_{\|}-$проекция импульса частицы на направление $\boldsymbol{H}$, a интеграл берется вдоль траектории между точками поворота частицы.





\end{document}